\section{A2: Occupation numbers and Fock states}

Let $\{\phi_j\}$ be an orthonormal one-particle basis. If $n_j$ bosons occupy
orbital $j$, then
\[
\sum_j n_j=N,
\qquad
\boxed{|n_0,n_1,n_2,\ldots\rangle}
\]
denotes the corresponding occupation-number or Fock state.

For example, two bosons in $a$ and one in $b$ have wavefunction
\[
\Psi_{2,1}
=\frac1{\sqrt3}(aab+aba+baa)
\]
and occupation representation $|2_a,1_b\rangle$. The latter records the
physical information without assigning identities to the bosons.

For occupations $\{n_j\}$, the number of distinct assignments to $N$
coordinate slots is
\[
M=\frac{N!}{\prod_j n_j!}.
\]
Hence the normalized symmetric wavefunction is
\[
\boxed{
\Psi_{\{n_j\}}
=
\sqrt{\frac{\prod_jn_j!}{N!}}
\sum_{\text{distinct }P}
\phi_{\alpha_{P(1)}}(\mathbf r_1)\cdots
\phi_{\alpha_{P(N)}}(\mathbf r_N),
}
\]
where orbital $j$ appears $n_j$ times in the list
$(\alpha_1,\ldots,\alpha_N)$.

If all bosons occupy $\phi_0$, then $n_0=N$, $M=1$, and
\[
\Psi_N=\prod_{i=1}^N\phi_0(\mathbf r_i),
\qquad
|\Psi_N\rangle=|N,0,0,\ldots\rangle.
\]

Occupation states are orthonormal:
\[
\langle m_0,m_1,\ldots|n_0,n_1,\ldots\rangle
=\prod_j\delta_{m_jn_j}.
\]
A general fixed-$N$ state is therefore
\[
|\Psi_N\rangle
=\sum_{\{n_j\};\,\sum_jn_j=N}
C_{n_0,n_1,\ldots}|n_0,n_1,\ldots\rangle,
\qquad
\sum_{\{n_j\}}|C_{n_0,n_1,\ldots}|^2=1.
\]

The bosonic Fock space collects all particle-number sectors:
\[
\mathcal F_B
=\mathcal H_0\oplus\mathcal H_1\oplus
\mathcal H_2^{(B)}\oplus\cdots.
\]
This formal enlargement does not imply particle-number nonconservation.

A simple condensate has one macroscopically occupied orbital:
\[
n_0=O(N),
\qquad
\lim_{N\to\infty}\frac{n_0}{N}>0.
\]
Thus condensation requires a finite fraction of all bosons in one orbital,
rather than literally every particle.
